\appendix
\chapter{C标准库函数速查手册}

本章汇总 C语言标准库中最常用的函数，按头文件分类，提供函数原型、功能说明和简要示例。

\section{C标准库头文件总览}

\begin{table}[H]
\centering
\small
\renewcommand{\arraystretch}{1.5}
\begin{tabulary}{\linewidth}{CC}
\toprule
\textbf{头文件} & \textbf{主要功能} \\ \midrule
\texttt{<stdio.h>} & 标准输入输出、文件操作、格式化 I/O \\ 
\texttt{<stdlib.h>} & 内存管理、数值转换、随机数、程序控制、排序查找 \\ 
\texttt{<string.h>} & 字符串操作、内存操作 \\ 
\texttt{<math.h>} & 数学函数（三角、指数、对数、取整等） \\ 
\texttt{<ctype.h>} & 字符分类与大小写转换 \\ 
\texttt{<time.h>} & 时间与日期处理 \\ 
\texttt{<assert.h>} & 运行时断言调试 \\ 
\texttt{<errno.h>} & 错误码定义 \\ 
\texttt{<stdarg.h>} & 可变参数处理 \\ 
\texttt{<limits.h>} & 整型类型取值范围常量 \\ 
\texttt{<float.h>} & 浮点型精度与范围常量 \\ 
\texttt{<stddef.h>} & 常用类型定义（size\_t, ptrdiff\_t, NULL 等） \\ 
\texttt{<setjmp.h>} & 非局部跳转（setjmp/longjmp） \\ 
\texttt{<signal.h>} & 信号处理 \\ 
\texttt{<locale.h>} & 本地化设置 \\ 
\texttt{<stdbool.h>} & 布尔类型支持（C99） \\ 
\texttt{<inttypes.h>} & 固定宽度整型与格式化宏（C99） \\ \bottomrule
\end{tabulary}
\caption{C标准库头文件一览}
\label{tab:stdlib_overview}
\end{table}

\section{\texorpdfstring{\texttt{<stdio.h>} 标准输入输出库}{<stdio.h> 标准输入输出库}}

\subsection{格式化输入输出}

\begin{table}[H]
\centering
\renewcommand{\arraystretch}{1.6}
\begin{tabulary}{\linewidth}{CCC}
\toprule
\textbf{函数原型} & \textbf{功能说明} & \textbf{返回值} \\ \midrule
\texttt{int printf(const char* fmt, ...)} & 格式化输出到 stdout & 成功输出的字符数 \\ 
\texttt{int scanf(const char* fmt, ...)} & 从 stdin 格式化输入 & 成功赋值的参数个数 \\ 
\texttt{int fprintf(FILE* fp, const char* fmt, ...)} & 格式化输出到文件 & 成功输出的字符数 \\ 
\texttt{int fscanf(FILE* fp, const char* fmt, ...)} & 从文件格式化输入 & 成功赋值的参数个数 \\ 
\texttt{int sprintf(char* buf, const char* fmt, ...)} & 格式化输出到字符串缓冲区 & 写入的字符数（不含 \textbackslash 0） \\ 
\texttt{int snprintf(char* buf, size\_t n, const char* fmt, ...)} & 安全版 sprintf，限制最大写入 n 字节 & 应写入的字符数 \\ 
\texttt{int sscanf(const char* str, const char* fmt, ...)} & 从字符串格式化输入 & 成功赋值的参数个数 \\ \bottomrule
\end{tabulary}
\caption{stdio.h 格式化输入输出函数}
\label{tab:stdio_format}
\end{table}

\subsection{字符与字符串 I/O}

\begin{table}[H]
\centering
\small
\renewcommand{\arraystretch}{1.6}
\begin{tabulary}{\linewidth}{CCC}
\toprule
\textbf{函数原型} & \textbf{功能说明} & \textbf{返回值} \\ \midrule
\texttt{int getchar(void)} & 从 stdin 读取一个字符 & 读取的字符（EOF 表示结束） \\ 
\texttt{int putchar(int c)} & 向 stdout 输出一个字符 & 输出的字符（EOF 表示失败） \\ 
\texttt{int fgetc(FILE* fp)} & 从文件流读取一个字符 & 读取的字符（EOF 表示结束） \\ 
\texttt{int fputc(int c, FILE* fp)} & 向文件流写入一个字符 & 写入的字符（EOF 表示失败） \\ 
\texttt{char* fgets(char* s, int n, FILE* fp)} & 从文件流读取最多 n-1 个字符 & s 指针（NULL 表示失败/结束） \\ 
\texttt{int fputs(const char* s, FILE* fp)} & 向文件流写入字符串（不含 \textbackslash 0） & 非负值成功，EOF 失败 \\ 
\texttt{char* gets(char* s)} & \textbf{已废弃}：从 stdin 读取一行，无长度限制 & s 指针（\textbf{存在缓冲区溢出风险}） \\ \bottomrule
\end{tabulary}
\caption{stdio.h 字符与字符串 I/O 函数}
\label{tab:stdio_char}
\end{table}

\subsection{文件操作}

\begin{table}[H]
\centering
\small
\renewcommand{\arraystretch}{1.5}
\begin{tabulary}{\linewidth}{CCC}
\toprule
\textbf{函数原型} & \textbf{功能说明} & \textbf{返回值} \\ \midrule
\texttt{FILE* fopen(const char* path, const char* mode)} & 以指定模式打开文件 & FILE 指针（NULL 表示失败） \\ 
\texttt{int fclose(FILE* fp)} & 关闭文件并刷新缓冲区 & 0 成功，EOF 失败 \\ 
\texttt{size\_t fread(void* ptr, size\_t sz, size\_t n, FILE* fp)} & 从文件读取 n 个大小为 sz 的数据块 & 实际读取的块数 \\ 
\texttt{size\_t fwrite(const void* ptr, size\_t sz, size\_t n, FILE* fp)} & 向文件写入 n 个大小为 sz 的数据块 & 实际写入的块数 \\ 
\texttt{int fseek(FILE* fp, long off, int whence)} & 移动文件位置指针 & 0 成功，非零失败 \\ 
\texttt{long ftell(FILE* fp)} & 获取当前文件位置偏移量 & 偏移量（-1L 表示失败） \\ 
\texttt{void rewind(FILE* fp)} & 将文件位置指针重置到开头 & 无 \\ 
\texttt{int feof(FILE* fp)} & 检测是否到达文件末尾 & 非零表示已到末尾 \\ 
\texttt{int ferror(FILE* fp)} & 检测文件流是否发生错误 & 非零表示有错误 \\ 
\texttt{int remove(const char* path)} & 删除指定文件 & 0 成功，非零失败 \\ 
\texttt{int rename(const char* old, const char* new)} & 重命名文件 & 0 成功，非零失败 \\ 
\texttt{void clearerr(FILE* fp)} & 清除文件流的错误和 EOF 标志 & 无 \\ 
\texttt{void setbuf(FILE* fp, char* buf)} & 设置文件流的缓冲区 & 无 \\ 
\texttt{int fflush(FILE* fp)} & 刷新输出缓冲区，将数据写入文件 & 0 成功，EOF 失败 \\ \bottomrule
\end{tabulary}
\caption{stdio.h 文件操作函数}
\label{tab:stdio_file}
\end{table}

\subsection{其他 I/O 与错误处理}

\begin{table}[H]
\centering
\small
\renewcommand{\arraystretch}{1.5}
\begin{tabulary}{\linewidth}{CCC}
\toprule
\textbf{函数原型} & \textbf{功能说明} & \textbf{返回值} \\ \midrule
\texttt{int ungetc(int c, FILE* fp)} & 将字符 c 推回流中 & 推入的字符（EOF 表示失败） \\ 
\texttt{void perror(const char* msg)} & 将错误信息输出到 stderr & 格式为 \texttt{msg: strerror(errno)} \\ 
\texttt{int setvbuf(FILE* fp, char* buf, int mode, size\_t size)} & 设置流的缓冲模式与缓冲区 & 0 成功，非零失败 \\ 
\texttt{FILE* tmpfile(void)} & 创建临时二进制文件（关闭自动删除） & FILE 指针（NULL 表示失败） \\ 
\texttt{char* tmpnam(char* s)} & 生成唯一的临时文件名 & 指向文件名的指针 \\ \bottomrule
\end{tabulary}
\caption{stdio.h 其他 I/O 与错误处理函数}
\label{tab:stdio_other}
\end{table}

\section{\texorpdfstring{\texttt{<stdlib.h>} 标准通用工具库}{<stdlib.h> 标准通用工具库}}

\subsection{内存管理}

\begin{table}[H]
\centering
\small
\renewcommand{\arraystretch}{1.6}
\begin{tabulary}{\linewidth}{CCC}
\toprule
\textbf{函数原型} & \textbf{功能说明} & \textbf{返回值} \\ \midrule
\texttt{void* malloc(size\_t size)} & 分配 size 字节的未初始化内存 & 指向分配内存的指针（NULL 失败） \\ 
\texttt{void* calloc(size\_t n, size\_t sz)} & 分配 n*sz 字节并初始化为 0 & 指向分配内存的指针（NULL 失败） \\ 
\texttt{void* realloc(void* ptr, size\_t sz)} & 重新调整已分配内存的大小 & 新指针（NULL 失败，原内存不变） \\ 
\texttt{void free(void* ptr)} & 释放之前分配的内存 & 无 \\ \bottomrule
\end{tabulary}
\caption{stdlib.h 内存管理函数}
\label{tab:stdlib_mem}
\end{table}

\subsection{数值转换}

\begin{table}[H]
\centering
\small
\renewcommand{\arraystretch}{1.6}
\begin{tabulary}{\linewidth}{CCC}
\toprule
\textbf{函数原型} & \textbf{功能说明} & \textbf{返回值} \\ \midrule
\texttt{int atoi(const char* str)} & 字符串转 int & 转换后的整数值 \\ 
\texttt{long atol(const char* str)} & 字符串转 long & 转换后的长整数值 \\ 
\texttt{long long atoll(const char* str)} & 字符串转 long long & 转换后的 long long 值 \\ 
\texttt{double atof(const char* str)} & 字符串转 double & 转换后的浮点值 \\ 
\texttt{long strtol(const char* s, char** end, int base)} & 字符串转 long，支持指定进制 & 转换结果，end 指向首个无效字符 \\ 
\texttt{double strtod(const char* s, char** end)} & 字符串转 double & 转换结果，end 指向首个无效字符 \\ \bottomrule
\end{tabulary}
\caption{stdlib.h 数值转换函数}
\label{tab:stdlib_conv}
\end{table}

\subsection{随机数与程序控制}

\begin{table}[H]
\centering
\small
\renewcommand{\arraystretch}{1.6}
\begin{tabulary}{\linewidth}{CCC}
\toprule
\textbf{函数原型} & \textbf{功能说明} & \textbf{返回值} \\ \midrule
\texttt{int rand(void)} & 生成 $[0, \text{RAND\_MAX}]$ 范围内的伪随机数 & 随机整数值 \\ 
\texttt{void srand(unsigned int seed)} & 设置随机数种子 & 无 \\ 
\texttt{void exit(int status)} & 立即终止程序，刷新并关闭所有流 & 不返回 \\ 
\texttt{void abort(void)} & 异常终止程序，生成 SIGABRT 信号 & 不返回 \\ 
\texttt{int atexit(void (*func)(void))} & 注册程序正常终止时调用的函数 & 0 成功，非零失败 \\ 
\texttt{int system(const char* cmd)} & 执行系统命令 & 命令返回值 \\ \bottomrule
\end{tabulary}
\caption{stdlib.h 随机数与程序控制函数}
\label{tab:stdlib_ctrl}
\end{table}

\subsection{查找与排序}

\begin{table}[H]
\centering
\small
\renewcommand{\arraystretch}{1.6}
\begin{tabulary}{\linewidth}{CCC}
\toprule
\textbf{函数原型} & \textbf{功能说明} & \textbf{返回值} \\ \midrule
\texttt{void qsort(void* base, size\_t n, size\_t sz, int (*cmp)(const void*, const void*))} & 对数组进行快速排序 & 无（原地排序） \\ 
\texttt{void* bsearch(const void* key, const void* base, size\_t n, size\_t sz, int (*cmp)(const void*, const void*))} & 在已排序数组中二分查找 & 指向匹配元素的指针（NULL 未找到） \\ \bottomrule
\end{tabulary}
\caption{stdlib.h 查找与排序函数}
\label{tab:stdlib_search}
\end{table}

\subsection{绝对值与除法}

\begin{table}[H]
\centering
\small
\renewcommand{\arraystretch}{1.5}
\begin{tabulary}{\linewidth}{CCC}
\toprule
\textbf{函数原型} & \textbf{功能说明} & \textbf{返回值} \\ \midrule
\texttt{int abs(int n)} & 计算整数的绝对值 & 绝对值 \\ 
\texttt{long labs(long n)} & 计算 long 的绝对值 & 绝对值 \\ 
\texttt{long long llabs(long long n)} & 计算 long long 的绝对值 (C99) & 绝对值 \\ 
\texttt{div\_t div(int num, int den)} & 整数除法（返回商和余数） & \texttt{div\_t} 结构体 \{quot, rem\} \\ 
\texttt{ldiv\_t ldiv(long num, long den)} & long 整数除法 & \texttt{ldiv\_t} 结构体 \\ \bottomrule
\end{tabulary}
\caption{stdlib.h 绝对值与除法函数}
\label{tab:stdlib_math}
\end{table}

\subsection{多字节与宽字符转换}

\begin{table}[H]
\centering
\small
\renewcommand{\arraystretch}{1.5}
\begin{tabulary}{\linewidth}{CCC}
\toprule
\textbf{函数原型} & \textbf{功能说明} & \textbf{返回值} \\ \midrule
\texttt{int mblen(const char* s, size\_t n)} & 获取多字节字符的字节数 & 字节数（-1 表示无效） \\ 
\texttt{int mbtowc(wchar\_t* wc, const char* s, size\_t n)} & 多字节字符转宽字符 & 字节数 \\ 
\texttt{int wctomb(char* s, wchar\_t wc)} & 宽字符转多字节字符 & 字节数 \\ 
\texttt{size\_t mbstowcs(wchar\_t* dest, const char* src, size\_t n)} & 多字节字符串转宽字符串 & 转换的宽字符数 \\ 
\texttt{size\_t wcstombs(char* dest, const wchar\_t* src, size\_t n)} & 宽字符串转多字节字符串 & 转换的字节数 \\ \bottomrule
\end{tabulary}
\caption{stdlib.h 多字节与宽字符转换函数}
\label{tab:stdlib_mb}
\end{table}

\section{\texorpdfstring{\texttt{<string.h>} 字符串处理库}{<string.h> 字符串处理库}}

\subsection{字符串操作}

\begin{table}[H]
\centering
\small
\renewcommand{\arraystretch}{1.6}
\begin{tabulary}{\linewidth}{CCC}
\toprule
\textbf{函数原型} & \textbf{功能说明} & \textbf{返回值} \\ \midrule
\texttt{size\_t strlen(const char* s)} & 计算字符串长度（不含 \textbackslash 0） & 字符个数 \\ 
\texttt{char* strcpy(char* dst, const char* src)} & 复制字符串（含 \textbackslash 0） & dst 指针 \\ 
\texttt{char* strncpy(char* dst, const char* src, size\_t n)} & 最多复制 n 个字符 & dst 指针 \\ 
\texttt{char* strcat(char* dst, const char* src)} & 拼接字符串到 dst 末尾 & dst 指针 \\ 
\texttt{char* strncat(char* dst, const char* src, size\_t n)} & 最多拼接 n 个字符 & dst 指针 \\ 
\texttt{int strcmp(const char* s1, const char* s2)} & 字典序比较两个字符串 & $<0$ s1小，$=0$ 相等，$>0$ s1大 \\ 
\texttt{int strncmp(const char* s1, const char* s2, size\_t n)} & 最多比较前 n 个字符 & 同上 \\ 
\texttt{char* strchr(const char* s, int c)} & 查找字符 c 首次出现的位置 & 指向该位置的指针（NULL 未找到） \\ 
\texttt{char* strrchr(const char* s, int c)} & 查找字符 c 最后一次出现的位置 & 指向该位置的指针（NULL 未找到） \\ 
\texttt{char* strstr(const char* haystack, const char* needle)} & 查找子串首次出现的位置 & 指向子串起始位置的指针（NULL 未找到） \\ 
\texttt{char* strtok(char* str, const char* delim)} & 按分隔符拆分字符串（内部维护状态） & 下一个 token 指针（NULL 表示结束） \\ 
\texttt{char* strdup(const char* s)} & 复制字符串（内部调用 malloc） & 新分配的字符串指针 \\ \bottomrule
\end{tabulary}
\caption{string.h 字符串操作函数}
\label{tab:string_ops}
\end{table}

\subsection{内存操作}

\begin{table}[H]
\centering
\small
\renewcommand{\arraystretch}{1.6}
\begin{tabulary}{\linewidth}{CCC}
\toprule
\textbf{函数原型} & \textbf{功能说明} & \textbf{返回值} \\ \midrule
\texttt{void* memcpy(void* dst, const void* src, size\_t n)} & 复制 n 字节（\textbf{不可重叠}） & dst 指针 \\ 
\texttt{void* memmove(void* dst, const void* src, size\_t n)} & 复制 n 字节（\textbf{可安全重叠}） & dst 指针 \\ 
\texttt{void* memset(void* ptr, int val, size\_t n)} & 将前 n 字节设置为 val & ptr 指针 \\ 
\texttt{int memcmp(const void* s1, const void* s2, size\_t n)} & 比较前 n 个字节 & $<0, =0, >0$ \\ 
\texttt{void* memchr(const void* s, int c, size\_t n)} & 在前 n 字节中查找值 c & 指向该位置的指针（NULL 未找到） \\ \bottomrule
\end{tabulary}
\caption{string.h 内存操作函数}
\label{tab:string_mem}
\end{table}

\subsection{其他字符串操作}

\begin{table}[H]
\centering
\small
\renewcommand{\arraystretch}{1.5}
\begin{tabulary}{\linewidth}{CCC}
\toprule
\textbf{函数原型} & \textbf{功能说明} & \textbf{返回值} \\ \midrule
\texttt{char* strerror(int errnum)} & 获取错误码对应的错误描述字符串 & 指向错误信息的指针 \\ 
\texttt{size\_t strspn(const char* s, const char* accept)} & 计算 s 开头完全由 accept 中字符组成的长度 & 匹配的字符数 \\ 
\texttt{size\_t strcspn(const char* s, const char* reject)} & 计算 s 开头不包含 reject 中字符的长度 & 匹配的字符数 \\ 
\texttt{char* strpbrk(const char* s, const char* accept)} & 查找 s 中首次出现 accept 中任意字符的位置 & 指向该位置的指针（NULL 未找到） \\ 
\texttt{int strcoll(const char* s1, const char* s2)} & 根据当前 locale 比较字符串 & $<0, =0, >0$ \\ 
\texttt{size\_t strxfrm(char* dst, const char* src, size\_t n)} & 将字符串转换为 locale 相关的比较格式 & 转换后的长度 \\ \bottomrule
\end{tabulary}
\caption{string.h 其他字符串操作函数}
\label{tab:string_other}
\end{table}

\section{\texorpdfstring{\texttt{<math.h>} 数学函数库}{<math.h> 数学函数库}}

\subsection{常用数学函数}

\begin{table}[H]
\centering
\small
\renewcommand{\arraystretch}{1.6}
\begin{tabulary}{\linewidth}{CCC}
\toprule
\textbf{函数原型} & \textbf{功能说明} & \textbf{返回值} \\ \midrule
\texttt{double sqrt(double x)} & 平方根 $\sqrt{x}$ & 平方根值 \\ 
\texttt{double pow(double x, double y)} & 幂运算 $x^y$ & 幂值 \\ 
\texttt{double exp(double x)} & 自然指数 $e^x$ & 指数值 \\ 
\texttt{double log(double x)} & 自然对数 $\ln x$ & 对数值 \\ 
\texttt{double log10(double x)} & 常用对数 $\log_{10} x$ & 对数值 \\ 
\texttt{double fabs(double x)} & 绝对值 $|x|$ & 绝对值 \\ 
\texttt{double ceil(double x)} & 向上取整 $\lceil x \rceil$ & 不小于 x 的最小整数 \\ 
\texttt{double floor(double x)} & 向下取整 $\lfloor x \rfloor$ & 不大于 x 的最大整数 \\ 
\texttt{double round(double x)} & 四舍五入 & 最接近的整数 \\ 
\texttt{double fmod(double x, double y)} & 浮点取余 $x - n \cdot y$ & 余数 \\ \bottomrule
\end{tabulary}
\caption{math.h 常用数学函数}
\label{tab:math_funcs}
\end{table}

\subsection{三角函数}

\begin{table}[H]
\centering
\small
\renewcommand{\arraystretch}{1.6}
\begin{tabulary}{\linewidth}{CCC}
\toprule
\textbf{函数原型} & \textbf{功能说明} & \textbf{返回值} \\ \midrule
\texttt{double sin(double x)} & 正弦（x 为弧度） & $\sin x$ \\ 
\texttt{double cos(double x)} & 余弦（x 为弧度） & $\cos x$ \\ 
\texttt{double tan(double x)} & 正切（x 为弧度） & $\tan x$ \\ 
\texttt{double asin(double x)} & 反正弦，定义域 $[-1, 1]$ & 弧度值 $[-\pi/2, \pi/2]$ \\ 
\texttt{double acos(double x)} & 反余弦，定义域 $[-1, 1]$ & 弧度值 $[0, \pi]$ \\ 
\texttt{double atan(double x)} & 反正切 & 弧度值 $(-\pi/2, \pi/2)$ \\ 
\texttt{double atan2(double y, double x)} & 双参数反正切（判定象限） & 弧度值 $(-\pi, \pi]$ \\ \bottomrule
\end{tabulary}
\caption{math.h 三角函数}
\label{tab:math_trig}
\end{table}

\subsection{双曲函数}

\begin{table}[H]
\centering
\small
\renewcommand{\arraystretch}{1.5}
\begin{tabulary}{\linewidth}{CCC}
\toprule
\textbf{函数原型} & \textbf{功能说明} & \textbf{返回值} \\ \midrule
\texttt{double sinh(double x)} & 双曲正弦 $\frac{e^x - e^{-x}}{2}$ & 双曲正弦值 \\ 
\texttt{double cosh(double x)} & 双曲余弦 $\frac{e^x + e^{-x}}{2}$ & 双曲余弦值 \\ 
\texttt{double tanh(double x)} & 双曲正切 $\frac{\sinh x}{\cosh x}$ & 双曲正切值 \\ \bottomrule
\end{tabulary}
\caption{math.h 双曲函数}
\label{tab:math_hyper}
\end{table}

\subsection{其他数学函数}

\begin{table}[H]
\centering
\small
\renewcommand{\arraystretch}{1.5}
\begin{tabulary}{\linewidth}{CCC}
\toprule
\textbf{函数原型} & \textbf{功能说明} & \textbf{返回值} \\ \midrule
\texttt{double cbrt(double x)} & 立方根 $\sqrt[3]{x}$ (C99) & 立方根值 \\ 
\texttt{double hypot(double x, double y)} & 计算 $\sqrt{x^2 + y^2}$（避免溢出） & 斜边长度 \\ 
\texttt{double log2(double x)} & 以 2 为底的对数 (C99) & 对数值 \\ 
\texttt{double trunc(double x)} & 截断小数部分 (C99) & 整数部分 \\ 
\texttt{double fmax(double x, double y)} & 返回两数中的较大值 (C99) & 较大值 \\ 
\texttt{double fmin(double x, double y)} & 返回两数中的较小值 (C99) & 较小值 \\ 
\texttt{int isnan(double x)} & 判断是否为 NaN (C99) & 非零为真，零为假 \\ 
\texttt{int isinf(double x)} & 判断是否为无穷大 (C99) & 非零为真，零为假 \\ 
\texttt{int isfinite(double x)} & 判断是否为有限值 (C99) & 非零为真，零为假 \\ \bottomrule
\end{tabulary}
\caption{math.h 其他数学函数}
\label{tab:math_other}
\end{table}

\vspace{0.5em}
\noindent\textbf{常量}：\texttt{M\_PI}（$\pi \approx 3.14159$）、\texttt{M\_E}（$e \approx 2.71828$）。编译时需链接数学库：\texttt{gcc -lm}。

\section{\texorpdfstring{\texttt{<ctype.h>} 字符分类与转换库}{<ctype.h> 字符分类与转换库}}

\begin{table}[H]
\centering
\renewcommand{\arraystretch}{1.6}
\begin{tabulary}{\linewidth}{CCC}
\toprule
\textbf{函数原型} & \textbf{功能说明} & \textbf{返回值} \\ \midrule
\texttt{int isalpha(int c)} & 是否为字母（A-Z, a-z） & 非零为真，零为假 \\ 
\texttt{int isdigit(int c)} & 是否为数字（0-9） & 非零为真，零为假 \\ 
\texttt{int isalnum(int c)} & 是否为字母或数字 & 非零为真，零为假 \\ 
\texttt{int isspace(int c)} & 是否为空白字符（空格、\textbackslash t、\textbackslash n 等） & 非零为真，零为假 \\ 
\texttt{int isupper(int c)} & 是否为大写字母 & 非零为真，零为假 \\ 
\texttt{int islower(int c)} & 是否为小写字母 & 非零为真，零为假 \\ 
\texttt{int ispunct(int c)} & 是否为标点符号 & 非零为真，零为假 \\ 
\texttt{int iscntrl(int c)} & 是否为控制字符 & 非零为真，零为假 \\ 
\texttt{int isgraph(int c)} & 是否为图形字符（除空格外的可打印字符） & 非零为真，零为假 \\ 
\texttt{int isxdigit(int c)} & 是否为十六进制数字（0-9, a-f, A-F） & 非零为真，零为假 \\ 
\texttt{int isblank(int c)} & 是否为空白字符（空格或制表符，C99） & 非零为真，零为假 \\ 
\texttt{int isprint(int c)} & 是否为可打印字符 & 非零为真，零为假 \\ 
\texttt{int toupper(int c)} & 转为大写 & 大写字符 \\ 
\texttt{int tolower(int c)} & 转为小写 & 小写字符 \\ \bottomrule
\end{tabulary}
\caption{ctype.h 字符分类与转换函数}
\label{tab:ctype_funcs}
\end{table}

\section{\texorpdfstring{\texttt{<time.h>} 时间与日期库}{<time.h> 时间与日期库}}

\begin{table}[H]
\centering
\renewcommand{\arraystretch}{1.6}
\begin{tabulary}{\linewidth}{CCC}
\toprule
\textbf{函数原型} & \textbf{功能说明} & \textbf{返回值} \\ \midrule
\texttt{time\_t time(time\_t* t)} & 获取当前日历时间（秒数） & 从 1970-01-01 起的秒数 \\ 
\texttt{double difftime(time\_t t2, time\_t t1)} & 计算两个时间之差 & 秒数差值 \\ 
\texttt{struct tm* localtime(const time\_t* t)} & 将时间戳转为本地时间结构体 & tm 结构体指针 \\ 
\texttt{struct tm* gmtime(const time\_t* t)} & 将时间戳转为 UTC 时间结构体 & tm 结构体指针 \\ 
\texttt{time\_t mktime(struct tm* tm)} & 将 tm 结构体转为时间戳 & time\_t 值 \\ 
\texttt{char* asctime(const struct tm* tm)} & 将 tm 结构体转为可读字符串 & 字符串指针 \\ 
\texttt{char* ctime(const time\_t* t)} & 将时间戳转为可读字符串 & 字符串指针 \\ 
\texttt{size\_t strftime(char* s, size\_t max, const char* fmt, const struct tm* tm)} & 按格式将 tm 转为字符串 & 写入的字符数 \\ 
\texttt{clock\_t clock(void)} & 获取程序运行至今的 CPU 时间 & clock\_t 值（除以 CLOCKS\_PER\_SEC 得秒） \\ \bottomrule
\end{tabulary}
\caption{time.h 时间与日期函数}
\label{tab:time_funcs}
\end{table}

\noindent\begin{minipage}{\textwidth}
\begin{lstlisting}[caption={time.h 使用示例}]
#include <stdio.h>
#include <time.h>

int main() {
    time_t now = time(NULL);
    struct tm* local = localtime(&now);
    printf("当前时间: %d-%02d-%02d %02d:%02d:%02d\n",
           local->tm_year + 1900, local->tm_mon + 1, local->tm_mday,
           local->tm_hour, local->tm_min, local->tm_sec);
    return 0;
}
\end{lstlisting}
\end{minipage}

\vspace{0.5em}
\begin{sloppypar}
\noindent\textbf{结构体说明}：\texttt{struct timespec} 包含 \texttt{tv\_sec}（秒数）和 \texttt{tv\_nsec}（纳秒数），用于高精度时间表示（POSIX/C11）。
\end{sloppypar}

\section{\texorpdfstring{\texttt{<assert.h>} 断言库}{<assert.h> 断言库}}

\begin{table}[H]
\centering
\renewcommand{\arraystretch}{1.6}
\begin{tabulary}{\linewidth}{CCC}
\toprule
\textbf{宏/函数} & \textbf{功能说明} & \textbf{行为} \\ \midrule
\texttt{assert(expression)} & 运行时断言检查 & 表达式为假时打印错误信息并调用 abort() \\ 
\texttt{NDEBUG} & 宏定义 & 在 \texttt{\#include <assert.h>} 前定义 \texttt{NDEBUG} 可禁用所有 assert \\ \bottomrule
\end{tabulary}
\caption{assert.h 断言宏}
\label{tab:assert_funcs}
\end{table}

\noindent\begin{minipage}{\textwidth}
\begin{lstlisting}[caption={assert 使用示例}]
#include <assert.h>

int divide(int a, int b) {
    assert(b != 0);  // 如果 b 为 0，程序终止并输出错误信息
    return a / b;
}
\end{lstlisting}
\end{minipage}

\section{\texorpdfstring{\texttt{<errno.h>} 错误码库}{<errno.h> 错误码库}}

\begin{table}[H]
\centering
\renewcommand{\arraystretch}{1.6}
\begin{tabulary}{\linewidth}{CCC}
\toprule
\textbf{宏/变量} & \textbf{功能说明} & \textbf{典型值} \\ \midrule
\texttt{errno} & 全局错误码变量 & 由库函数在出错时设置 \\ 
\texttt{EDOM} & 数学函数定义域错误 & 如 sqrt(-1) \\ 
\texttt{ERANGE} & 结果超出表示范围 & 如溢出 \\ 
\texttt{EAGAIN} & 资源暂时不可用 & 非阻塞 I/O 场景 \\ 
\texttt{EINVAL} & 无效参数 & 传入非法参数 \\ \bottomrule
\end{tabulary}
\caption{errno.h 错误码}
\label{tab:errno_funcs}
\end{table}

\noindent\begin{minipage}{\textwidth}
\begin{lstlisting}[caption={errno 使用示例}]
#include <stdio.h>
#include <errno.h>
#include <string.h>
#include <math.h>

int main() {
    errno = 0;
    double result = sqrt(-1.0);
    if (errno == EDOM) {
        printf("数学定义域错误: %s\n", strerror(errno));
    }
    return 0;
}
\end{lstlisting}
\end{minipage}

\section{\texorpdfstring{\texttt{<stdarg.h>} 可变参数库}{<stdarg.h> 可变参数库}}

\begin{table}[H]
\centering
\renewcommand{\arraystretch}{1.6}
\begin{tabulary}{\linewidth}{CCC}
\toprule
\textbf{宏/类型} & \textbf{功能说明} & \textbf{用法} \\ \midrule
\texttt{va\_list} & 可变参数列表类型 & 用于声明参数遍历变量 \\ 
\texttt{va\_start(ap, last)} & 初始化 va\_list & last 是最后一个固定参数名 \\ 
\texttt{va\_arg(ap, type)} & 获取下一个参数 & type 为期望的参数类型 \\ 
\texttt{va\_end(ap)} & 清理 va\_list & 在函数返回前调用 \\ \bottomrule
\end{tabulary}
\caption{stdarg.h 可变参数宏}
\label{tab:stdarg_funcs}
\end{table}

\noindent\begin{minipage}{\textwidth}
\begin{lstlisting}[caption={可变参数函数示例}]
#include <stdio.h>
#include <stdarg.h>

double average(int count, ...) {
    va_list args;
    va_start(args, count);
    double sum = 0.0;
    for (int i = 0; i < count; i++) {
        sum += va_arg(args, double);
    }
    va_end(args);
    return sum / count;
}

int main() {
    printf("平均值: %.2f\n", average(4, 10.0, 20.0, 30.0, 40.0));
    return 0;
}
\end{lstlisting}
\end{minipage}

\section{\texorpdfstring{\texttt{<limits.h>} 与 \texttt{<float.h>} 极限值常量}{<limits.h> 与 <float.h> 极限值常量}}

\subsection{limits.h 整型极限}

\begin{table}[H]
\centering
\renewcommand{\arraystretch}{1.5}
\begin{tabulary}{\linewidth}{CCC}
\toprule
\textbf{宏} & \textbf{含义} & \textbf{典型值} \\ \midrule
\texttt{CHAR\_BIT} & char 的位数 & 8 \\ 
\texttt{CHAR\_MIN / CHAR\_MAX} & char 最小/最大值 & -128 / 127 \\ 
\texttt{INT\_MIN / INT\_MAX} & int 最小/最大值 & -2147483648 / 2147483647 \\ 
\texttt{LONG\_MIN / LONG\_MAX} & long 最小/最大值 & 平台相关 \\ 
\texttt{LLONG\_MIN / LLONG\_MAX} & long long 最小/最大值 & $\approx \pm 9.22 \times 10^{18}$ \\ 
\texttt{UINT\_MAX} & unsigned int 最大值 & 4294967295 \\ \bottomrule
\end{tabulary}
\caption{limits.h 常用宏常量}
\label{tab:limits_consts}
\end{table}

\subsection{float.h 浮点型极限}

\begin{table}[H]
\centering
\renewcommand{\arraystretch}{1.5}
\begin{tabulary}{\linewidth}{CCC}
\toprule
\textbf{宏} & \textbf{含义} & \textbf{典型值} \\ \midrule
\texttt{FLT\_MIN / DBL\_MIN} & float/double 最小正值 & $1.175 \times 10^{-38}$ / $2.225 \times 10^{-308}$ \\ 
\texttt{FLT\_MAX / DBL\_MAX} & float/double 最大值 & $3.402 \times 10^{38}$ / $1.797 \times 10^{308}$ \\ 
\texttt{FLT\_EPSILON / DBL\_EPSILON} & 机器精度 $\epsilon$ & $1.192 \times 10^{-7}$ / $2.220 \times 10^{-16}$ \\ 
\texttt{FLT\_DIG / DBL\_DIG} & 有效十进制位数 & 6 / 15 \\ \bottomrule
\end{tabulary}
\caption{float.h 常用宏常量}
\label{tab:float_consts}
\end{table}

\section{\texorpdfstring{\texttt{<stddef.h>} 标准定义库}{<stddef.h> 标准定义库}}

\begin{table}[H]
\centering
\small
\renewcommand{\arraystretch}{1.5}
\begin{tabulary}{\linewidth}{CCC}
\toprule
\textbf{宏/类型} & \textbf{说明} & \textbf{典型值/用法} \\ \midrule
\texttt{size\_t} & 无符号整型，用于表示对象大小 & \texttt{sizeof} 的返回类型 \\ 
\texttt{ptrdiff\_t} & 有符号整型，用于表示指针差值 & 两个指针相减的结果类型 \\ 
\texttt{wchar\_t} & 宽字符类型，用于存储宽字符 & 通常为 2 或 4 字节 \\ 
\texttt{NULL} & 空指针常量 & 通常定义为 \texttt{((void*)0)} 或 \texttt{0} \\ 
\texttt{offsetof(type, member)} & 计算结构体成员相对于起始地址的偏移量 & \texttt{offsetof(struct S, a)} \\ \bottomrule
\end{tabulary}
\caption{stddef.h 标准定义}
\label{tab:stddef}
\end{table}

\section{\texorpdfstring{\texttt{<setjmp.h>} 非局部跳转库}{<setjmp.h> 非局部跳转库}}

\begin{table}[H]
\centering
\small
\renewcommand{\arraystretch}{1.5}
\begin{tabulary}{\linewidth}{CCC}
\toprule
\textbf{宏/类型} & \textbf{功能说明} & \textbf{返回值/行为} \\ \midrule
\texttt{jmp\_buf} & 数组类型，用于保存调用环境（寄存器状态、栈指针等） & 供 \texttt{setjmp} 和 \texttt{longjmp} 使用 \\ 
\texttt{int setjmp(jmp\_buf env)} & 保存当前执行环境到 \texttt{env} 中 & 直接调用返回 0；通过 \texttt{longjmp} 返回非零值 \\ 
\texttt{void longjmp(jmp\_buf env, int val)} & 恢复到 \texttt{env} 保存的环境，并继续从 \texttt{setjmp} 处执行 & \texttt{setjmp} 将返回 \texttt{val}（若 val 为 0 则返回 1） \\ \bottomrule
\end{tabulary}
\caption{setjmp.h 非局部跳转宏}
\label{tab:setjmp}
\end{table}

\section{\texorpdfstring{\texttt{<signal.h>} 信号处理库}{<signal.h> 信号处理库}}

\begin{table}[H]
\centering
\small
\renewcommand{\arraystretch}{1.5}
\begin{tabulary}{\linewidth}{CCC}
\toprule
\textbf{宏/函数} & \textbf{功能说明} & \textbf{典型值/行为} \\ \midrule
\texttt{SIGINT} & 交互式注意信号（通常由 Ctrl+C 触发） & 值为 2 \\ 
\texttt{SIGTERM} & 终止请求信号 & 值为 15 \\ 
\texttt{SIGSEGV} & 无效内存访问（段错误） & 值为 11 \\ 
\texttt{SIG\_DFL} & 默认信号处理动作 & 恢复系统默认行为 \\ 
\texttt{SIG\_IGN} & 忽略信号 & 信号被丢弃，不执行任何操作 \\ 
\texttt{signal(int sig, void (*func)(int))} & 注册信号处理函数 & 返回旧的处理函数指针 \\ 
\texttt{raise(int sig)} & 向当前进程发送指定信号 & 成功返回 0，失败返回非零 \\ \bottomrule
\end{tabulary}
\caption{signal.h 信号处理宏与函数}
\label{tab:signal}
\end{table}

\section{\texorpdfstring{\texttt{<locale.h>} 本地化设置库}{<locale.h> 本地化设置库}}

\begin{table}[H]
\centering
\small
\renewcommand{\arraystretch}{1.5}
\begin{tabulary}{\linewidth}{CCC}
\toprule
\textbf{宏/函数} & \textbf{功能说明} & \textbf{典型值/行为} \\ \midrule
\texttt{LC\_ALL} & 设置所有本地化类别 & 涵盖数字、货币、时间、字符等 \\ 
\texttt{LC\_COLLATE} & 影响字符串比较和排序规则 & 如 \texttt{strcoll} 的行为 \\ 
\texttt{LC\_CTYPE} & 影响字符分类与转换函数 & 如 \texttt{isalpha}, \texttt{toupper} \\ 
\texttt{LC\_MONETARY} & 影响货币格式化 & 如 \texttt{localeconv} 返回的货币符号 \\ 
\texttt{LC\_NUMERIC} & 影响数字格式化（如小数点符号） & 如 \texttt{printf} 的小数点 \\ 
\texttt{LC\_TIME} & 影响日期和时间格式化 & 如 \texttt{strftime} 的输出 \\ 
\texttt{char* setlocale(int category, const char* locale)} & 设置或查询当前程序的本地化环境 & \texttt{""} 表示使用系统环境变量 \\ 
\texttt{struct lconv* localeconv(void)} & 获取当前本地化环境的数字和货币格式信息 & 返回指向 \texttt{lconv} 结构体的指针 \\ \bottomrule
\end{tabulary}
\caption{locale.h 本地化设置宏与函数}
\label{tab:locale}
\end{table}

\section{\texorpdfstring{\texttt{<stdbool.h>} 布尔类型库 (C99)}{<stdbool.h> 布尔类型库 (C99)}}

\begin{table}[H]
\centering
\small
\renewcommand{\arraystretch}{1.5}
\begin{tabulary}{\linewidth}{CC}
\toprule
\textbf{宏/类型} & \textbf{说明} \\ \midrule
\texttt{bool} & 扩展为 \texttt{\_Bool} 类型 \\ 
\texttt{true} & 扩展为整数常量 1 \\ 
\texttt{false} & 扩展为整数常量 0 \\ 
\texttt{\_\_bool\_true\_false\_are\_defined} & 扩展为 1，表示宏已定义 \\ \bottomrule
\end{tabulary}
\caption{stdbool.h 布尔类型宏}
\label{tab:stdbool}
\end{table}

\section{\texorpdfstring{\texttt{<inttypes.h>} 固定宽度整型库 (C99)}{<inttypes.h> 固定宽度整型库 (C99)}}

\begin{table}[H]
\centering
\small
\renewcommand{\arraystretch}{1.5}
\begin{tabulary}{\linewidth}{CCC}
\toprule
\textbf{类型/宏} & \textbf{说明} & \textbf{典型值} \\ \midrule
\texttt{int8\_t / uint8\_t} & 精确 8 位有符号/无符号整数 & 范围 $\pm 127$ / $0 \sim 255$ \\ 
\texttt{int16\_t / uint16\_t} & 精确 16 位有符号/无符号整数 & 范围 $\pm 32767$ / $0 \sim 65535$ \\ 
\texttt{int32\_t / uint32\_t} & 精确 32 位有符号/无符号整数 & 范围 $\pm 2\times 10^9$ / $0 \sim 4\times 10^9$ \\ 
\texttt{int64\_t / uint64\_t} & 精确 64 位有符号/无符号整数 & 范围 $\pm 9\times 10^{18}$ \\ 
\texttt{PRIu32 / PRId64} & \texttt{printf} 格式化固定宽度整型的宏 & 如 \texttt{PRIu32} 展开为 \texttt{"u"} \\ 
\texttt{SCNd32 / SCNu64} & \texttt{scanf} 格式化固定宽度整型的宏 & 如 \texttt{SCNd32} 展开为 \texttt{"d"} \\ \bottomrule
\end{tabulary}
\caption{inttypes.h 固定宽度整型与格式化宏}
\label{tab:inttypes}
\end{table}

\section{\texorpdfstring{\texttt{<tgmath.h>} 泛型数学宏库 (C99)}{<tgmath.h> 泛型数学宏库 (C99)}}

\begin{table}[H]
\centering
\small
\renewcommand{\arraystretch}{1.5}
\begin{tabulary}{\linewidth}{CCC}
\toprule
\textbf{泛型宏} & \textbf{功能说明} & \textbf{行为} \\ \midrule
\texttt{sin(x), cos(x), tan(x)} & 根据参数类型自动选择 float/double/long double 版本 & 传入 float 调用 \texttt{sinf}，传入 double 调用 \texttt{sin} \\ 
\texttt{sqrt(x), pow(x, y)} & 自动匹配数学函数变体 & 避免手动类型转换 \\ 
\texttt{fabs(x), ceil(x), floor(x)} & 自动选择绝对值、取整函数 & 支持整型与浮点型参数 \\ 
\texttt{atan2(y, x), hypot(x, y)} & 自动选择双参数函数版本 & 参数类型不同时提升为精度更高者 \\ \bottomrule
\end{tabulary}
\caption{tgmath.h 泛型数学宏}
\label{tab:tgmath}
\end{table}

\section{\texorpdfstring{\texttt{<wchar.h>} 宽字符 I/O 库 (C95)}{<wchar.h> 宽字符 I/O 库 (C95)}}

\begin{table}[H]
\centering
\small
\renewcommand{\arraystretch}{1.5}
\begin{tabulary}{\linewidth}{CCC}
\toprule
\textbf{函数原型} & \textbf{功能说明} & \textbf{返回值} \\ \midrule
\texttt{wint\_t fgetwc(FILE* fp)} & 从文件流读取一个宽字符 & 读取的宽字符（WEOF 表示结束） \\ 
\texttt{wint\_t fputwc(wchar\_t wc, FILE* fp)} & 向文件流写入一个宽字符 & 写入的宽字符（WEOF 表示失败） \\ 
\texttt{wchar\_t* fgetws(wchar\_t* s, int n, FILE* fp)} & 从文件流读取宽字符串 & s 指针（NULL 表示失败） \\ 
\texttt{int fputws(const wchar\_t* s, FILE* fp)} & 向文件流写入宽字符串 & 非负值成功，EOF 失败 \\ 
\texttt{wint\_t getwchar(void)} & 从 stdin 读取一个宽字符 & 读取的宽字符 \\ 
\texttt{wint\_t putwchar(wchar\_t wc)} & 向 stdout 输出一个宽字符 & 输出的宽字符 \\ 
\texttt{size\_t wcslen(const wchar\_t* s)} & 计算宽字符串长度 & 宽字符个数 \\ 
\texttt{wchar\_t* wcscpy(wchar\_t* dst, const wchar\_t* src)} & 复制宽字符串 & dst 指针 \\ 
\texttt{int wcscmp(const wchar\_t* s1, const wchar\_t* s2)} & 比较宽字符串 & $<0, =0, >0$ \\ 
\texttt{wint\_t btowc(int c)} & 单字节字符转宽字符 & 宽字符值 \\ 
\texttt{int wctob(wint\_t wc)} & 宽字符转单字节字符 & 单字节值（EOF 表示无法转换） \\ \bottomrule
\end{tabulary}
\caption{wchar.h 宽字符 I/O 与操作函数}
\label{tab:wchar}
\end{table}

\section{\texorpdfstring{\texttt{<wctype.h>} 宽字符分类库 (C95)}{<wctype.h> 宽字符分类库 (C95)}}

\begin{table}[H]
\centering
\small
\renewcommand{\arraystretch}{1.5}
\begin{tabulary}{\linewidth}{CCC}
\toprule
\textbf{函数原型} & \textbf{功能说明} & \textbf{返回值} \\ \midrule
\texttt{int iswalnum(wint\_t wc)} & 是否为宽字母或数字 & 非零为真，零为假 \\ 
\texttt{int iswalpha(wint\_t wc)} & 是否为宽字母 & 非零为真，零为假 \\ 
\texttt{int iswdigit(wint\_t wc)} & 是否为宽数字 & 非零为真，零为假 \\ 
\texttt{int iswspace(wint\_t wc)} & 是否为宽空白字符 & 非零为真，零为假 \\ 
\texttt{int iswlower(wint\_t wc)} & 是否为宽小写字母 & 非零为真，零为假 \\ 
\texttt{int iswupper(wint\_t wc)} & 是否为宽大写字母 & 非零为真，零为假 \\ 
\texttt{wint\_t towlower(wint\_t wc)} & 宽字符转小写 & 小写宽字符 \\ 
\texttt{wint\_t towupper(wint\_t wc)} & 宽字符转大写 & 大写宽字符 \\ 
\texttt{wctype\_t wctype(const char* property)} & 获取字符属性描述符 & 属性类型 \\ 
\texttt{int iswctype(wint\_t wc, wctype\_t desc)} & 根据描述符检查宽字符属性 & 非零为真，零为假 \\ \bottomrule
\end{tabulary}
\caption{wctype.h 宽字符分类与转换函数}
\label{tab:wctype}
\end{table}

\section{\texorpdfstring{\texttt{<complex.h>} 复数运算库 (C99)}{<complex.h> 复数运算库 (C99)}}

\begin{table}[H]
\centering
\small
\renewcommand{\arraystretch}{1.5}
\begin{tabulary}{\linewidth}{CCC}
\toprule
\textbf{宏/函数} & \textbf{功能说明} & \textbf{返回值} \\ \midrule
\texttt{complex / \_Complex} & 复数类型关键字 & 支持 float/double/long double \\ 
\texttt{I} & 虚数单位 $\sqrt{-1}$ & 复数常量 \\ 
\texttt{double creal(double complex z)} & 获取复数的实部 & 实部值 \\ 
\texttt{double cimag(double complex z)} & 获取复数的虚部 & 虚部值 \\ 
\texttt{double cabs(double complex z)} & 计算复数的模 $|z|$ & 模值 \\ 
\texttt{double carg(double complex z)} & 计算复数的辐角（相位） & 弧度值 $(-\pi, \pi]$ \\ 
\texttt{double complex conj(double complex z)} & 计算复共轭 & 共轭复数 \\ 
\texttt{double complex csqrt(double complex z)} & 复数平方根 & 平方根复数 \\ 
\texttt{double complex cexp(double complex z)} & 复数指数 $e^z$ & 复数值 \\ 
\texttt{double complex clog(double complex z)} & 复数自然对数 & 复数值 \\ 
\texttt{double complex cpow(double complex x, double complex y)} & 复数幂运算 $x^y$ & 复数值 \\ \bottomrule
\end{tabulary}
\caption{complex.h 复数运算函数}
\label{tab:complex}
\end{table}

\section{\texorpdfstring{\texttt{<fenv.h>} 浮点环境库 (C99)}{<fenv.h> 浮点环境库 (C99)}}

\begin{table}[H]
\centering
\small
\renewcommand{\arraystretch}{1.5}
\begin{tabulary}{\linewidth}{CCC}
\toprule
\textbf{宏/函数} & \textbf{功能说明} & \textbf{返回值/行为} \\ \midrule
\texttt{FE\_DIVBYZERO} & 除零异常标志 & 浮点异常类型 \\ 
\texttt{FE\_INEXACT} & 结果不精确异常标志 & 如 $\pi$ 的近似值 \\ 
\texttt{FE\_INVALID} & 无效操作异常标志 & 如 $\sqrt{-1}$ \\ 
\texttt{FE\_OVERFLOW} & 上溢异常标志 & 结果超出最大值 \\ 
\texttt{FE\_UNDERFLOW} & 下溢异常标志 & 结果低于最小正值 \\ 
\texttt{int feclearexcept(int excepts)} & 清除指定的浮点异常标志 & 0 成功 \\ 
\texttt{int fetestexcept(int excepts)} & 测试哪些异常标志被设置 & 被设置的标志位掩码 \\ 
\texttt{int feraiseexcept(int excepts)} & 手动触发浮点异常 & 0 成功 \\ 
\texttt{int fegetround(void)} & 获取当前浮点舍入模式 & 舍入模式常量 \\ 
\texttt{int fesetround(int round)} & 设置浮点舍入模式 & 0 成功，非零失败 \\ 
\texttt{fenv\_t} & 保存完整浮点环境的状态类型 & 供 \texttt{fegetenv}/\texttt{fesetenv} 使用 \\ \bottomrule
\end{tabulary}
\caption{fenv.h 浮点环境控制函数}
\label{tab:fenv}
\end{table}

\section{\texorpdfstring{\texttt{<iso646.h>} 替代拼写宏库 (C95)}{<iso646.h> 替代拼写宏库 (C95)}}

\begin{table}[H]
\centering
\small
\renewcommand{\arraystretch}{1.5}
\begin{tabulary}{\linewidth}{CC}
\toprule
\textbf{宏} & \textbf{替代运算符} \\ \midrule
\texttt{and} & \texttt{\&\&} \\ 
\texttt{or} & \texttt{||} \\ 
\texttt{not} & \texttt{!} \\ 
\texttt{bitand} & \texttt{\&} \\ 
\texttt{bitor} & \texttt{|} \\ 
\texttt{xor} & \texttt{\^} \\ 
\texttt{compl} & \texttt{\textasciitilde} \\ 
\texttt{and\_eq} & \texttt{\&=} \\ 
\texttt{or\_eq} & \texttt{|=} \\ 
\texttt{xor\_eq} & \texttt{\^=} \\ 
\texttt{not\_eq} & \texttt{!=} \\ \bottomrule
\end{tabulary}
\caption{iso646.h 替代拼写宏}
\label{tab:iso646}
\end{table}
